\documentclass{article}
\usepackage{titling}
\usepackage{sectsty}
\usepackage{siunitx}
\usepackage[dvipsnames]{xcolor}
\usepackage{fancyhdr}
\usepackage{lastpage}
\usepackage[colorlinks=true]{hyperref}
\hypersetup{
    linkcolor=black,     
    citecolor=MidnightBlue,
    urlcolor=MidnightBlue   
}
\usepackage{amssymb}
\usepackage{amsmath}
\usepackage{soul}

\pagestyle{fancy}
\fancyhf{}
\fancyfoot[C]{{\Large\textbf{\thepage}}/\pageref{LastPage}}
\renewcommand{\headrulewidth}{0pt}

\sectionfont{\color{blue!90!black}}
\newcommand{\mathSI}[2]{\text{\SI{#1}{#2}}}
\setulcolor{cyan}

\begin{document}

\begin{center}
	\textcolor{cyan}{\rule{0.75\textwidth}{0.6pt}}\\
    \vspace{1cm}
    {\Huge \bfseries \textcolor{blue!70!black}{International Astronomy and Astrophysics Competition}}\\
    \vspace{1cm}
    {\Large \bfseries \textcolor{black}{Qualification Round 2026}}\\
    \vspace{1cm}
    {\large \bfseries \textcolor{NavyBlue}{Solution Document}}\\
    \vspace{1.5cm}
    {\color{Cyan} \Huge $\diamond$}\\
    \vspace{1.5cm}
    {\Large \bfseries \textcolor{blue!90!black}{MD MOKHDUM AMIN FAHIM}}\\
    \vspace{0.5cm}
    {\Large Police Line Secondary School, Jashore}\\
    \vspace{0.5cm}
    {\Large \bfseries \textcolor{NavyBlue}{Bangladesh}}\\
    \vspace{1.5cm}
    {\color{Cyan} \Huge $\diamond$}\\
    \vspace{1.5cm}
    {\large Submission Date : \today }\\
    \vspace{0.5cm}
    {\large Number of Pages : \pageref{LastPage} }\\
    \vspace{1cm}
    \textcolor{Cyan}{\rule{0.75\textwidth}{0.6pt}}\\
\end{center}

\thispagestyle{empty}
\newpage
\setcounter{page}{1}

\section*{Problem A : Cosmic Distances}
\vspace{-10pt} \textcolor{cyan}{\rule{7.7cm}{0.6pt}}\\
\begin{enumerate}
    \item Pleiades (C)
    \item Large Magellanic Cloud (F)
    \item Vega (A)
    \item Coma Cluster (H)
    \item Sagittarius A* (E)
    \item Aldebaran (B)
    \item Eagle Nebula (D)
    \item Andromeda Galaxy (G)\\
\end{enumerate}

\section*{Problem B : Low Earth Orbit}
\vspace{-10pt} \textcolor{cyan}{\rule{7.5cm}{0.6pt}}\\\\
Given :
\begin{itemize}
    \item Radius of Earth $R=$ \mathSI{6370}{km}
    \item Above Earth’s surface, LEO starts a $r_1=$ \mathSI{160}{km}
    \item Above Earth’s surface, LEO ends at $r_2=$ \mathSI{2000}{km}
    \item Approximate number of satellites in LEO $N=13500$\\
\end{itemize}
{\Large \bfseries \textcolor{blue!90!black}{\ul{(a)}}}\\\\
Volume of sphere or radius $r_1$ : $V_1=\frac{4}{3}\pi (R+r_1)^3$\\
Volume of sphere or radius $r_2$ : $V_1=\frac{4}{3}\pi (R+r_2)^3$\\\\
$\therefore$ Volume of space within LEO :
$$V=V_2-V_1=\frac{4}{3}\pi (R+r_2)^3-\frac{4}{3}\pi (R+r_1)^3=\frac{4}{3}\pi \left[(R+r_2)^3-(R+r_1)^3\right]$$
$$V=\frac{4}{3}\pi\left[(6370 \text{ km}+ 2000\text{ km})^3-(6370\text{ km}+ 160\text{ km})^3\right]$$
\[\therefore V \approx 1.29\times10^{12} \text{ km}^3\]\\
\[
\boxed{\textcolor{blue!90!black}{\textbf{Ans : }} V \approx 1.29\times10^{12} \text{ km}^3}
\]

\newpage
\noindent
{\Large \bfseries \textcolor{blue!90!black}{\ul{(b)}}}\\\\
There are approximately $N=13500$ satellites in LEO which has space approximately $V = 1.29\times10^{12} \text{ km}^3 $\\\\
$\therefore$  The average volume of space available to each satellite :
\[\langle V\rangle=\frac{V}{N}=\frac{1.29\times10^{12} \text{ km}^3}{13500} \]
\[\therefore \langle V\rangle \approx 9.56\times10^7\text{ km}^3\]\\
\[
\boxed{\textcolor{blue!90!black}{\textbf{Ans : }} \langle V\rangle \approx 9.56\times10^7 \text{ km}^3}
\]\\
\section*{Problem C : Total Solar Eclipse}
\vspace{-10pt} \textcolor{cyan}{\rule{8cm}{0.6pt}}\\\\
{\Large \bfseries \textcolor{blue!90!black}{\ul{(a)}}}\\\\
 Total, Partial, and Annular eclipses are three types of Solar eclipses. A solar eclipse is an astronomical phenomenon which occurs when the Moon is positioned between the Sun and Earth, consequently casting a shadow on Earth as the Moon blocks the Sun's light. The main difference of the three types of Solar eclipses is about how and how much Sun's light is blocked.\\\\
\textbf{Total Solar Eclipse :} It occurs when the Moon is directly between the Earth and the Sun and the Moon's angular diameter is equal or larger than the Sun's. The Moon blocks the Sun's light totally, just revealing the solar corona. \\\\
\textbf{Partial Solar Eclipse :} It occurs when the Sun, Moon, and Earth are not perfectly in a line. So, the Moon covers a portion of the Sun, resulting the Sun to look like a crescent.  \\\\
\textbf{Annular Solar Eclipse :} This happens when the Moon's angular diameter is smaller than the Sun's and the Moon is at it's apogee. As the Moon appears smaller than the Sun, it cannot fully cover it. As a result Sun's light makes a bright ring around the Moon's shadow.\\\\
\newpage
\noindent
{\Large \bfseries \textcolor{blue!90!black}{\ul{(b)}}}\\\\
According to \href{https://science.nasa.gov/eclipses/future-eclipses/}{NASA Science}, the next total solar eclipse will occur in 12 August, 2026 and it is after the Qualification Round deadline, which is in April 2026. A total solar eclipse will be visible in Greenland, Iceland, Spain, Russia, and a small area of Portugal, while a partial eclipse will be visible in Europe, Africa, North America, the Atlantic Ocean, Arctic Ocean, and Pacific Ocean. \\
\[
\boxed{\textcolor{blue!90!black}{\textbf{Ans : }} 12~August,~2026}
\]\\\\
{\Large \bfseries \textcolor{blue!90!black}{\ul{(c)}}}\\\\
Given,
\begin{itemize}
    \item  The radii of the Moon and the Sun are $R_M$ and $R_S$
    \item The Earth‑Moon and Earth‑Sun distances are $d_M$ and $d_S$\\
\end{itemize}
The case is approximated by taking the Moon's orbit circular and the diameter of Moon as an arc on that circle. Therefore, the angular diameter of the Moon is,
$$\theta_M=\frac{2R_M}{d_M}$$\\
It is similarly approximated in the case of the Sun. Therefore, the angular diameter of the Moon is,
$$\theta_S=\frac{2R_S}{d_S}$$\\
Total solar eclipse occurs when the Moon's angular diameter is equal or larger than the Sun's. Therefore,
$$\theta_M \geq \theta_S$$
By substitution, we get :
$$\therefore\frac{2R_M}{d_M}\geq \frac{2R_S}{d_S} $$
\[
\boxed{\therefore \frac{R_M}{d_M}\geq \frac{R_S}{d_S} }
\]
Thus, the required relation is proved.
\newpage
\section*{Problem D : The Last Eclipse}
\vspace{-10pt} \textcolor{cyan}{\rule{7.6cm}{0.6pt}}\\\\
Given,
\begin{itemize}
    \item The Moon's radius $R_M=1740\text{ km}$
    \item The Sun's present radius $R_S(0)=R_{S_0}=696000\text{ km}$ 
    \item The Earth-Sun aphelion distance $d_S=152000000\text{ km}$
    \item Present Earth-Moon perigee distance $d_{M_0}=356000\text{ km}$
    \item Rate of distance of the Moon's moving away $v=3.8\text{ cm/yr}=3.8\times10^{-5}\text{ km/yr}$
    \item According to Gough(1981), the increase in the Sun’s radius can be approximated by,
    \[\frac{R_S(t)}{R_S(0)}=\sqrt{\frac{1}{1-\frac{2}{5}\frac{t}{t_0}}}\]
    where $R_{S} (t)=R_{S}$ is the Sun’s radius at time $t$ from today, and\\ $t_0=4.6 \text{ billion years}=4.6\times10^9\text{ yr}$.\\
\end{itemize}
The Earth-Moon perigee distance after time $t$,
\[d_M=d_{M_0}+vt\]
The Sun’s radius after time $t$,
\[R_S=\frac{R_{S_0}}{\sqrt{1-\frac{2}{5}\frac{t}{t_0}}}\]
The angular diameter of the Moon is,
$$\theta_M=\frac{2R_M}{d_M}$$\\
The angular diameter of the Moon is,
$$\theta_S=\frac{2R_S}{d_S}$$\\
The Earth should be at aphelion and the Moon should be at perigee during the last total solar eclipse. The angular diameter of the Moon and the Sun should be equal as it is the last total solar eclipse. 
\newpage
Therefore,
\[\theta_M= \theta_S\]
\[\frac{2R_M}{d_M}=\frac{2R_S}{d_S}\]
\[\frac{R_M}{d_{M_0}+vt}=\frac{R_{S_0}}{d_S\sqrt{1-\frac{2}{5}\frac{t}{t_0}}}\]
Simple algebraic manipulation yields,
\[(5t_0v^2R_{S_0}^2)\cdot t^2 + (10t_0v d_{M_0}R_{S_0}^2+2d_S^2R_M^2)\cdot t + 5t_0(R_{S_0}^2d_{M_0}^2-R_M^2d_S^2)=0\]
Comparing the equation with $at^2+bt+c=0$, \\
$a=5t_0v^2R_{S_0}^2=1.608842419\times10^{13} \text{ km}^4/\text{yr}$\\
$b=10t_0v d_{M_0}R_{S_0}^2+2d_S^2R_M^2=4.413456046\times10^{23}\text{ km}^4 $\\
$c=5t_0(R_{S_0}^2d_{M_0}^2-R_M^2d_S^2)=-1.968046572\times10^{32}\text{ km}^4\text{yr}$\\
The solutions of the quadratic equation are,
\[t_{1,2}=\frac{-b±\sqrt{b^2-4ac}}{2a}\]
$t_1\approx438900000\text{ yr}\approx 439 \text{ Million yr}$\\
$t_2\approx-2.7871\times10^{10}\text{ yr}$\\
Negative value of time is not acceptable. So,
\[t=439 \text{ Million yr}\]\\
\[
\boxed{\textcolor{blue!90!black}{\textbf{Ans : }} 439 \text{ Million year}}
\] \\
\section*{Problem E : Testing the Theory of Relativity}
\vspace{-10pt} \textcolor{cyan}{\rule{11.5cm}{0.6pt}}\\\\
The total solar eclipse of 29 May 1919 was used to test Albert Einstein’s General Theory of Relativity through expeditions led by Sir Arthur Eddington to Príncipe and by Andrew Crommelin and Charles Davidson to Sobral, Brazil. Unlike Newton's Law of Gravitation, Einstein proposed that mass and energy curves spacetime, meaning starlight should be deflected as it passes the Sun. The teams observed a deflection of roughly 1.75 arcseconds that confirmed Einstein’s prediction.

\end{document}
